\documentclass[../../../include/open-logic-section]{subfiles}

\begin{document}

\olfileid{his}{set}{infinitesimals}
\olsection{لامتناہی طور اُتے چھوٹیاں قدراں تے تفریق}

نیوٹن تے لائبنز نے ستارویں صدی دے آخر وچ حسابان وکھو وکھ لبھیا۔ حسابان دی اک خاص اہم ورتوں \emph{تفریق} سی۔ موٹے طور اُتے تفریق دا مقصد ایہ دَسنا اے کہ کسے نقطے اُتے فنکشن دی «بدلن دی شرح»، یا ڈھلوان، کیہ اے۔

ایہ خیال اک صاف مثال نال سمجھ آ سکدا اے۔ تھلّے کالے رنگ وچ وکھایا فنکشن
$f(x) = \nicefrac{x^2}{4} + \nicefrac{1}{2}$ لو:
\begin{center}
	\begin{tikzpicture}[scale=1]
	\draw[->, gray] (-1,0) -- (4.25,0) node[right] {$x$};
	\draw[->, gray] (0,-0.25) -- (0,5.25) node[above] {$f(x)$};
	
	\foreach \x/\xtext in {1/1, 2/2, 3/3, 4/4}
	\draw[shift={(\x,0)}, gray] (0pt,2pt) -- (0pt,-2pt) node[below] {$\xtext$};
	
	\foreach \y/\ytext in {1/1, 2/2, 3/3, 4/4, 5/5}
	\draw[shift={(0,\y)}, gray] (2pt,0pt) -- (-2pt,0pt) node[left] {$\ytext$};
	
	\draw[oldiagcolorC] (0.5, 9/16) -- (3.5, 57/16) -- (3.5, 9/16)--cycle;
	\draw[oldiagcolorD] (.5, 9/16) -- (2.5, 33/16) -- (2.5, 9/16) -- cycle;
	\draw[oldiagcolorE] (.5, 9/16) -- (1.5, 17/16) -- (1.5, 9/16) -- cycle;
	\draw[thick] (-1,.75) parabola bend (0,0.5) (4,4.5);
	\end{tikzpicture}
\end{center}
فرض کرو اسیں $c =
\nicefrac{1}{2}$ اُتے فنکشن دی ڈھلوان لبھنی اے۔ اسیں پہلاں اک تکون بناؤنے آں جیہدے وتر دی ڈھلوان اوس نقطے اُتے فنکشن دی ڈھلوان دے نیڑے ہووے۔ اُتلی لال تکون لے لو۔ جے ساڈی تکون دے پاوے دی لمبائی $\beta$ اے، تاں اوہدی اُچائی
$f(\nicefrac{1}{2}+\beta) - f(\nicefrac{1}{2})$ اے، سو وتر دی ڈھلوان اے:
\[
\frac{f(\nicefrac{1}{2}+\beta) - f(\nicefrac{1}{2})}{\beta}.
\]
سو !!{colorC} تکون دے پاوے دی لمبائی~$3$ اے تے اوہدی ڈھلوان عین~$1$ اے۔ اوس توں چھوٹی !!{colorD} تکون دا پاوا~$2$ اے۔ اوہدا وتر چنگا اندازہ دیندا اے، اوہدی ڈھلوان $\nicefrac{3}{4}$ اے۔ ہور وی چھوٹی !!{colorE} تکون دا پاوا~$1$ اے۔ ایہدے نال اندازہ ہور وی نیڑے اے، تے ڈھلوان $\nicefrac{1}{2}$ اے۔

جِنّی تکون چھوٹی، اُنا ای اندازہ اصل دے نیڑے۔ سو اسیں ایویں آکھ سکدے آں کہ جیہڑی تکون دے پاوے دی لمبائی \emph{لامتناہی طور اُتے چھوٹی} ہووے، اوہدے وتر دی ڈھلوان $c =
\nicefrac{1}{2}$ اُتے عین فنکشن دی ڈھلوان اے۔ ایس طرح فنکشن $f$ دے (پہلے) مشتق لئی، نقطے $c$ اُتے، ایہ فارمولا ملدا اے:
\[
{f'}(c) = \frac{f(c+\beta) - f(c)}{\beta} \text{ جتھے $\beta$ لامتناہی طور اُتے چھوٹا اے۔}
\]
نیوٹن تے لائبنز نے موٹے طور اُتے ایہو گل آکھی۔

پر جدوں اوہناں ایہ گل آکھی، تاں اوہناں توں پچھنا پیندا اے کہ \emph{لامتناہی طور اُتے چھوٹی قدر} کیہ اے؟ ایتھے وڈی دُبدھا پیدا ہوندی اے۔ جے $\beta = 0$ ہووے، تاں $f'$ دی تعریف ٹھیک نہیں بندی، کیوں جے $0$ نال تقسیم کرنا پیندا اے۔ پر جے $\beta >
0$ ہووے، تاں ڈھلوان دا بس اک \emph{اندازہ} ملدا اے، عین ڈھلوان نہیں۔

ایہ فِکر بعدلے زمانے دی سوچ نوں پچھلے ویلے اُتے لاگو کر کے نہیں اُٹھدی۔ ویکھو برکلے نیوٹن دے پیروکاراں دی تنقید کِویں کردا اے:
\begin{quote}
میں منّدا آں کہ علامتاں کسے شے نوں وی دس سکدیاں نیں تے کجھ نہ ہون نوں وی۔ سو اصلی لکھت $c + \beta$ وچ $\beta$ اک وادھے یا کجھ نہ ہون دا نشان ہو سکدا سی۔ پر ایہناں وچوں جیہڑا وی مطلب تسی لوو، دلیل اوسے مطلب اُتے قائم رکھو۔ دُہرے مطلب نال اگے ودھنا صاف مغالطہ ہووے گا۔ (\citealt[\S{}XIII]{Berkeley1734}، متغیر پچھلی لکھت نال میل کرن لئی بدلے گئے نیں۔)
\end{quote}
برکلے دے اعتراض توں لامتناہی طور اُتے چھوٹیاں قدراں والے حسابان دا بچاؤ کردیاں کوئی آکھ سکدا اے کہ «لامتناہی طور اُتے چھوٹیاں قدراں» بس مجازی گل اے۔ کوئی ایہ وی آکھ سکدا اے کہ جد تک تکون \emph{سچ مُچ چھوٹی} لئے، سانوں مماس دا \emph{کم چلاؤن جوگا} اندازہ مل جاندا اے۔ برکلے دا ایہدا وی جواب سی۔ اوہ آکھدا سی کہ انجینئرنگ وچ ایہ کم چلاؤن جوگا ہو سکدا اے، پر ایہ ریاضی دی \emph{حیثیت} نوں کھورا لاؤندا اے، کیوں جے:
\begin{quote}
سانوں دَسیا جاندا اے کہ \emph{in rebus mathematicis errores qu\`{a}m minimi
non sunt contemnendi}. [ریاضی وچ نِکّی توں نِکّی غلطی وی نظرانداز نہیں کرنی چاہیدی۔] \citep[\S{}IX]{Berkeley1734}
\end{quote}
ترچھے حرفاں والا ٹوٹا نیوٹن دی اپنی کتاب \emph{Quadrature of Curves} (1704) توں لگ بھگ لفظ بہ لفظ لیا گیا اے۔

برکلے دے فلسفیانہ اعتراض ڈونگھائی نال گل دی جڑ تک جاندے نیں۔ فیر وی حسابان دی ورتوں بے حد کامیاب رہی تے ریاضی دان غلطی وچ پئے بغیر اوہنوں ورتدے رہے۔

\end{document}
